\documentclass[11pt,a4paper]{article}
% preamble.tex — estilo común de todos los apuntes Froth.
% Cada apunte hace % preamble.tex — estilo común de todos los apuntes Froth.
% Cada apunte hace % preamble.tex — estilo común de todos los apuntes Froth.
% Cada apunte hace \input{preamble} justo después de \documentclass.
\usepackage[margin=2.4cm]{geometry}
\usepackage{xcolor}
\definecolor{litblue}{RGB}{31,79,121}
\definecolor{litgreen}{RGB}{27,94,32}
\definecolor{litorange}{RGB}{230,81,0}
\usepackage{parskip}
\usepackage{titlesec}
\titleformat{\section}{\Large\bfseries\color{litblue}}{\thesection}{0.6em}{}
\titleformat{\subsection}{\large\bfseries\color{litblue}}{\thesubsection}{0.6em}{}
\usepackage{tcolorbox}
\usepackage{fancyhdr}
\pagestyle{fancy}
\fancyhf{}
\fancyhead[L]{\small\color{litblue}Froth — Apuntes de ML}
\fancyhead[R]{\small\thepage}
\renewcommand{\headrulewidth}{0.4pt}
\usepackage[colorlinks=true,linkcolor=litblue,urlcolor=litblue]{hyperref}

% Cabecera de cada apunte: \cabecera{numero}{titulo}
\newcommand{\cabecera}[2]{%
  \begin{tcolorbox}[colback=litblue,coltext=white,colframe=litblue,arc=2mm]
    {\Large\bfseries Apunte #1 — #2}\\[3pt]
    {\small Proyecto Froth · aprender Machine Learning construyendo}
  \end{tcolorbox}\vspace{0.8em}}

% Cajas temáticas
\newtcolorbox{concepto}[1]{colback=blue!4,colframe=litblue,title={Concepto: #1},arc=1.5mm}
\newtcolorbox{practica}{colback=green!5,colframe=litgreen,title={Pruébalo tú (teclado en mano)},arc=1.5mm}
\newtcolorbox{ojo}{colback=orange!8,colframe=litorange,title={Ojo / error típico},arc=1.5mm}
\newtcolorbox{resumen}{colback=gray!8,colframe=gray!60,title={En una frase},arc=1.5mm}

\newcommand{\codigo}[1]{\texttt{#1}}
 justo después de \documentclass.
\usepackage[margin=2.4cm]{geometry}
\usepackage{xcolor}
\definecolor{litblue}{RGB}{31,79,121}
\definecolor{litgreen}{RGB}{27,94,32}
\definecolor{litorange}{RGB}{230,81,0}
\usepackage{parskip}
\usepackage{titlesec}
\titleformat{\section}{\Large\bfseries\color{litblue}}{\thesection}{0.6em}{}
\titleformat{\subsection}{\large\bfseries\color{litblue}}{\thesubsection}{0.6em}{}
\usepackage{tcolorbox}
\usepackage{fancyhdr}
\pagestyle{fancy}
\fancyhf{}
\fancyhead[L]{\small\color{litblue}Froth — Apuntes de ML}
\fancyhead[R]{\small\thepage}
\renewcommand{\headrulewidth}{0.4pt}
\usepackage[colorlinks=true,linkcolor=litblue,urlcolor=litblue]{hyperref}

% Cabecera de cada apunte: \cabecera{numero}{titulo}
\newcommand{\cabecera}[2]{%
  \begin{tcolorbox}[colback=litblue,coltext=white,colframe=litblue,arc=2mm]
    {\Large\bfseries Apunte #1 — #2}\\[3pt]
    {\small Proyecto Froth · aprender Machine Learning construyendo}
  \end{tcolorbox}\vspace{0.8em}}

% Cajas temáticas
\newtcolorbox{concepto}[1]{colback=blue!4,colframe=litblue,title={Concepto: #1},arc=1.5mm}
\newtcolorbox{practica}{colback=green!5,colframe=litgreen,title={Pruébalo tú (teclado en mano)},arc=1.5mm}
\newtcolorbox{ojo}{colback=orange!8,colframe=litorange,title={Ojo / error típico},arc=1.5mm}
\newtcolorbox{resumen}{colback=gray!8,colframe=gray!60,title={En una frase},arc=1.5mm}

\newcommand{\codigo}[1]{\texttt{#1}}
 justo después de \documentclass.
\usepackage[margin=2.4cm]{geometry}
\usepackage{xcolor}
\definecolor{litblue}{RGB}{31,79,121}
\definecolor{litgreen}{RGB}{27,94,32}
\definecolor{litorange}{RGB}{230,81,0}
\usepackage{parskip}
\usepackage{titlesec}
\titleformat{\section}{\Large\bfseries\color{litblue}}{\thesection}{0.6em}{}
\titleformat{\subsection}{\large\bfseries\color{litblue}}{\thesubsection}{0.6em}{}
\usepackage{tcolorbox}
\usepackage{fancyhdr}
\pagestyle{fancy}
\fancyhf{}
\fancyhead[L]{\small\color{litblue}Froth — Apuntes de ML}
\fancyhead[R]{\small\thepage}
\renewcommand{\headrulewidth}{0.4pt}
\usepackage[colorlinks=true,linkcolor=litblue,urlcolor=litblue]{hyperref}

% Cabecera de cada apunte: \cabecera{numero}{titulo}
\newcommand{\cabecera}[2]{%
  \begin{tcolorbox}[colback=litblue,coltext=white,colframe=litblue,arc=2mm]
    {\Large\bfseries Apunte #1 — #2}\\[3pt]
    {\small Proyecto Froth · aprender Machine Learning construyendo}
  \end{tcolorbox}\vspace{0.8em}}

% Cajas temáticas
\newtcolorbox{concepto}[1]{colback=blue!4,colframe=litblue,title={Concepto: #1},arc=1.5mm}
\newtcolorbox{practica}{colback=green!5,colframe=litgreen,title={Pruébalo tú (teclado en mano)},arc=1.5mm}
\newtcolorbox{ojo}{colback=orange!8,colframe=litorange,title={Ojo / error típico},arc=1.5mm}
\newtcolorbox{resumen}{colback=gray!8,colframe=gray!60,title={En una frase},arc=1.5mm}

\newcommand{\codigo}[1]{\texttt{#1}}

\usepackage{tikz}
\begin{document}
\cabecera{05}{Qué es un embedding (el corazón del proyecto)}

\section{El problema: la computadora no entiende palabras}

Tú, como experto, sabes que estas dos frases hablan de lo mismo:
\begin{itemize}
  \item ``\emph{flotación de partículas finas}''
  \item ``\emph{recuperación de granos minerales pequeños mediante espuma}''
\end{itemize}
Pero no comparten casi ninguna palabra. Si buscáramos ``coincidencia de palabras'', el
sistema diría que no tienen nada que ver. Y al revés: ``banco de peces'' y ``banco de
dinero'' comparten la palabra ``banco'' y no se parecen en nada.

\textbf{Conclusión:} contar palabras iguales no captura el \emph{significado}. Necesitamos
otra cosa.

\section{La solución: convertir el texto en un punto en un mapa}

\begin{concepto}{embedding (vector de significado)}
Un \textbf{embedding} es un texto convertido en una lista de números —sus
\emph{coordenadas}— de modo que se puede colocar como un \textbf{punto} en un ``mapa del
significado''. La lista de números también se llama \textbf{vector}.

La propiedad mágica: el mapa está construido para que \textbf{textos con significado
parecido queden cerca}, y los distintos, lejos. Dos papers que estudian lo mismo caen
juntos aunque usen palabras diferentes.
\end{concepto}

Imagina el mapa de significado de tus papers (aquí en 2D para poder dibujarlo):

\begin{center}
\begin{tikzpicture}[scale=1.0]
  % zona flotación
  \fill[litblue!15] (1.2,1.2) ellipse (1.5 and 1.1);
  \node[litblue] at (1.2,2.6) {\small\textbf{flotación de finos}};
  \filldraw[litblue] (0.6,1.0) circle (2pt);
  \filldraw[litblue] (1.4,1.5) circle (2pt);
  \filldraw[litblue] (1.9,0.9) circle (2pt);
  \filldraw[litblue] (1.0,1.7) circle (2pt);
  % zona by-products
  \fill[litgreen!15] (6.0,1.4) ellipse (1.4 and 1.0);
  \node[litgreen] at (6.0,2.7) {\small\textbf{recuperación Ta/Nb/Be}};
  \filldraw[litgreen] (5.6,1.2) circle (2pt);
  \filldraw[litgreen] (6.3,1.6) circle (2pt);
  \filldraw[litgreen] (6.1,1.0) circle (2pt);
  % outlier off-topic
  \filldraw[litorange] (3.6,4.2) circle (2pt);
  \node[litorange] at (3.6,4.6) {\small paper de telecomunicaciones (ruido)};
  % ejes sin números
  \draw[->,gray] (-0.3,-0.3) -- (8,-0.3);
  \draw[->,gray] (-0.3,-0.3) -- (-0.3,5);
\end{tikzpicture}
\end{center}

Cada punto es un paper. Los de flotación se agrupan a la izquierda; los de by-products, a
la derecha; y el paper off-topic que se coló (¿recuerdas el de telecomunicaciones?) queda
\emph{solo, lejos de todo} --- por eso luego el sistema podrá apartarlo automáticamente.

\section{¿Por qué 768 coordenadas y no 2?}

\begin{concepto}{dimensiones}
En un mapa de verdad usas \textbf{2} coordenadas: latitud y longitud. Pero el significado
es mucho más rico que una posición en un plano, así que hace falta más ``espacio''. El
modelo que usaremos (SPECTER) describe cada texto con \textbf{$\sim$768 coordenadas}.

No lo puedes dibujar (nadie ve en 768 dimensiones), pero eso no importa: las
\textbf{matemáticas} sí saben medir qué tan cerca están dos puntos aunque tengan 768
coordenadas. En el Apunte 06 veremos cómo (la ``similitud coseno''). En la Fase 3
aplastaremos esas 768 a 2 justo para poder dibujar el mapa de verdad.
\end{concepto}

Cada uno de esos 768 números captura algún ``aspecto'' abstracto del texto. No son
interpretables uno a uno (el número 348 no significa ``cuánto habla de litio''); lo que
importa es el \emph{patrón completo} de los 768 juntos.

\section{¿De dónde salen los números? Un modelo preentrenado}

\begin{concepto}{modelo preentrenado}
Los 768 números los produce \textbf{SPECTER}: una \emph{red neuronal} que la organización
AllenAI ya entrenó leyendo millones de papers científicos (y quién cita a quién). Aprendió
a colocar textos parecidos cerca.

Lo clave: nosotros \textbf{no entrenamos nada}. Descargamos el modelo ya hecho y lo usamos,
como quien usa una calculadora sin fabricarla. Eso es un \emph{modelo preentrenado}, y es
la forma más rápida y potente de empezar en ML de texto.
\end{concepto}

\begin{ojo}
¿Por qué SPECTER y no un modelo genérico? Porque está entrenado \textbf{con papers}.
Un modelo generalista entiende texto común; SPECTER entiende el lenguaje académico (métodos,
resultados, jerga científica) mucho mejor para nuestro caso. Herramienta especializada para
un trabajo especializado.
\end{ojo}

\section{Qué haremos con esto (lo que viene)}

\begin{enumerate}
  \item Instalar las librerías que cargan SPECTER (siguiente paso).
  \item Convertir tus 126 abstracts en 126 vectores (una matriz $126 \times 768$).
  \item Medir parecidos con la similitud coseno (Apunte 06).
  \item Construir el \textbf{buscador semántico}: escribes una frase $\rightarrow$ te da los
        5 papers más parecidos por significado. (El momento ``wow''.)
\end{enumerate}

\begin{resumen}
Un embedding convierte cada texto en un punto (768 coordenadas) dentro de un ``mapa del
significado'' donde lo parecido queda cerca; esos números los da SPECTER, un modelo ya
entrenado con papers que nosotros solo reutilizamos.
\end{resumen}

\end{document}
